% ─────────────────────────────────────────────────────────────────────────────
% Checkliste "Anzeige der Vereinserrichtung" und nächste Schritte (Statuten-Generator,
% webapp: LatexStatutePdfRenderer). KEIN amtliches Formular -- die Anzeige selbst erfolgt
% formlos oder online bei der Vereinsbehörde (§ 11 VerG). RAW_-Variablen kommen escaped aus PHP.
% ─────────────────────────────────────────────────────────────────────────────
\documentclass[a4paper,10pt]{article}
\usepackage[utf8]{inputenc}
\usepackage[T1]{fontenc}
\usepackage[ngerman]{babel}
\usepackage[left=22mm,right=22mm,top=20mm,bottom=22mm]{geometry}
\usepackage{microtype}
\usepackage{parskip}
\usepackage{enumitem}
\usepackage{array}
\usepackage{tabularx}
\usepackage{amssymb}
\usepackage{fancyhdr}
\usepackage{lastpage}
\usepackage{xcolor}
\usepackage[scaled=0.96]{helvet}
\renewcommand{\familydefault}{\sfdefault}

\definecolor{navy}{RGB}{26,32,51}
\definecolor{midgray}{RGB}{107,114,128}
\definecolor{warnred}{RGB}{185,28,28}
\newcommand{\checkbox}{\ensuremath{\square}}
\newlist{checkliste}{itemize}{1}
\setlist[checkliste]{label=\checkbox,leftmargin=7mm,itemsep=1mm}

\def\entwurf{<<<RAW_ENTWURF>>>}
\pagestyle{fancy}
\fancyhf{}
\renewcommand{\headrulewidth}{0pt}
\ifx\entwurf\empty\else\fancyhead[C]{\small\bfseries\color{warnred}\entwurf}\fi
\fancyfoot[L]{\scriptsize\color{midgray}<<<VEREINSNAME>>> -- Checkliste zur Vereinsgründung, Stand <<<DATUM>>>}
\fancyfoot[R]{\scriptsize\color{midgray}Seite \thepage{} von \pageref{LastPage}}

\begin{document}
{\LARGE\bfseries\color{navy}Anzeige der Vereinserrichtung}\par
\vspace{1mm}
{\large\color{navy}Checkliste und nächste Schritte für \glqq{}<<<VEREINSNAME>>>\grqq{}}\par
\vspace{4mm}

\noindent{\renewcommand{\arraystretch}{1.3}%
\begin{tabularx}{\textwidth}{@{}>{\color{midgray}}p{48mm} X@{}}
Vereinsname & <<<VEREINSNAME>>> \\
Sitz & <<<SITZ>>> \\
Vereinsbehörde & <<<BEHOERDE>>> \\
Statuten & Fassung <<<VERSION>>> vom <<<DATUM>>> \\
Gründungsmitglieder & <<<GRUENDER_ANZAHL>>> \\
\end{tabularx}}

\section*{1. Vor der Anzeige}
\begin{checkliste}
  \item Statuten von einer fachkundigen Person prüfen lassen (dieses Dokument ist keine Rechtsberatung).
  \item Rechtsgrundlagen im RIS auf den aktuellen Stand prüfen (VerG, ElWG, EAG, UStG).
  \item Netzbetreiber kontaktieren: Lokal-/Regionalbereich und Zählpunkte bestätigen lassen.
  \item Gründungsversammlung abhalten und die Statuten beschließen (Protokollvorlage liegt bei).
\end{checkliste}

\section*{2. Anzeige an die Vereinsbehörde (§ 11 VerG)}
\begin{checkliste}
  \item Schriftlich oder online (oesterreich.gv.at) bei der Vereinsbehörde: <<<BEHOERDE>>>.
  \item Ein Exemplar der beschlossenen Statuten beilegen.
  \item Angaben zu jedem Gründungsmitglied: Name, Geburtsdatum, Geburtsort, Zustellanschrift
        (bei juristischen Personen: Name, Sitz, Registernummer, vertretungsbefugte Personen).
  \item Zustellanschrift des Vereins angeben.
  \item Die Behörde kann binnen vier Wochen (verlängerbar auf sechs Wochen) die Gründung untersagen oder zur
        Aufnahme der Tätigkeit einladen; erst danach entsteht der Verein (§ 13 VerG).
  \item Bekanntgabe der organschaftlichen Vertreter (Obmann/Obfrau, Kassier, ggf. weitere) mit Funktion, Name,
        Geburtsdatum, Geburtsort, Zustellanschrift und Beginn der Vertretungsbefugnis binnen vier Wochen nach der
        Bestellung.
\end{checkliste}
<<<RAW_GRUENDER>>>

\section*{3. Nach der Entstehung des Vereins}
\begin{checkliste}
  \item ZVR-Zahl notieren (Auszug aus dem Vereinsregister, kostenlos online abrufbar).
  \item Vereinskonto eröffnen; für SEPA-Lastschriften eine Creditor-ID bei der OeNB beantragen.
  \item Finanzamt: Steuernummer beantragen (Formular Verf 15a). Der Verein ist nicht gemeinnützig und mit seinen
        Gewinnen körperschaftsteuerpflichtig (keine Mindestkörperschaftsteuer).
  \item <<<RAW_UST>>>
  \item Vertrag mit dem Netzbetreiber über die Teilnahme an der gemeinsamen Energienutzung abschließen; Zählpunkte
        anmelden, Mitglieder erteilen die Datenfreigabe im Portal des Netzbetreibers.
  \item Mitgliedsverträge (Beitritt, Energiebezug, Einspeisung), Datenschutzinformation und ggf. AGB bereitstellen.
  \item Auftragsverarbeitungsverträge mit Dienstleistern abschließen (z. B. Abrechnungsplattform, Hosting).
  \item Prüfen, ob Lieferantenpflichten nach § 69 ElWG greifen (Anlagen über 100 kW der Gemeinschaft bzw.
        über 30 kW von Haushaltskunden).
  \item Einnahmen-Ausgaben-Rechnung samt Vermögensübersicht binnen fünf Monaten nach Ende jedes Rechnungsjahres
        (§ 21 VerG); Prüfung durch die Rechnungsprüfer binnen vier Monaten.
\end{checkliste}

\vspace{4mm}
{\small\color{midgray}Erstellt mit dem Statuten-Generator. Vorlage, keine Rechtsberatung. Vor Einreichung von einer
fachkundigen Person prüfen lassen.}
\end{document}
